\section{Localisation}
\label{section:high-level:localisation}

So far, we have assessed our code as one big black box.

\begin{itemize}
  \item We did not ask \emph{which part of the code} causes a problem. 
  \item We did not ask \emph{when} it arises throughout the execution.
  \item We did not ask \emph{where}, i.e.~on which rank (if we run over multiple
  nodes) a problem arises.
  \item We distinguished, to some degree, \emph{what} we want to study in our
  code, i.e.~which quantity of interest.
\end{itemize}

\noindent
Before we continue to assess the code and dive deeper into details, it makes sense to gather the information from the first three bullet points above or at least to discuss various approaches for gathering said information.
Hereby, we must consider the type of flaw we are investigating:
Single-core flaws imply that the ``where'' question does not play a significant
role.
Inter-node flaws might suggest that the ``where'' question is the primary
concern to address.

Finally, the most frequent quantity of interest (metric) will be time spent in routines.
In this case, the routines consuming most time are also called \emph{hotspots} of a code.
However, there are many other metrics that affect runtime, and we might want to study them separately—even though all of them eventually affect the time-to-solution\footnote{Some codes are interested in other metrics such as energy, but they are still rare.}.
Not all tools can provide answers for all metrics.
We therefore must carefully select our repertoire of tools.

\begin{instructions}{gprof}
t.b.d.
\end{instructions}


\begin{instructions}{Intel VTune}
t.b.d.
\end{instructions}

\paragraph*{How it works}

\begin{itemize}[noitemsep]
  \item[$\boxplus$] Once you have a first high-level overview, run a profiler to
  confirm that effect is either localised or uniformly found over all compute
  units.
  \item[$\boxplus$] If the flaw is found on few units only or only for a certain
  compute time or within one function which is not responsible for the majority
  of the runtime, study the flaw for this program phase, compute unit, function
  separatedly from all others.
  \item[$\boxplus$] Upon your very first profiling, ask a code expert which code parts (functions) actually perform calculations (cmp.~Section~\ref{section:preparation:labelling}).
\end{itemize}


\paragraph*{How it does not work}

\begin{itemize}[noitemsep]
  \item[$\boxminus$] Jump straight into some traces without spending time to
  think what metrics and program parts are relevant to answer your
  hypotheses regarding  performance flaws.
  \item[$\boxminus$] Draw conclusions from global observations on all code parts
  and compute resources (cores or ranks).
\end{itemize}
